\section{Cơ sở lý thuyết}
\label{sec:theory}

\subsection{Dự đoán RUL và các nghiên cứu liên quan}
Prognostics là lĩnh vực nghiên cứu diễn biến suy giảm của hệ thống và thời gian
còn lại trước khi hỏng. Saxena và cộng sự xây dựng C-MAPSS để có dữ liệu
run-to-failure với nhãn đã biết cho bài toán này~\cite{saxena2008damage}. Các
nghiên cứu sau đó tiếp cận theo hai hướng chính: trích xuất đặc trưng thủ công
rồi dùng mô hình hồi quy, hoặc học trực tiếp từ chuỗi bằng deep learning như
CNN và LSTM~\cite{lei2018machinery,li2018rul,zheng2017lstm}. Ramasso và Saxena
lưu ý rằng kết quả trên C-MAPSS phụ thuộc mạnh vào subset và cách đánh
giá~\cite{ramasso2014review}; vì vậy, các con số giữa những bài báo dùng cách
tiền xử lý hoặc chia dữ liệu khác nhau không so sánh trực tiếp được.

Báo cáo chọn hướng thứ nhất. Đặc trưng thống kê trên window kết hợp với mô hình
cây cho phép kiểm tra từng bước của pipeline và phù hợp với trọng tâm của môn
học là xử lý dữ liệu tái lập được. Bảng~\ref{tab:related-work} tóm tắt vị trí
của báo cáo so với các nguồn tham khảo chính.

\begin{table}[H]
\centering
\caption{Vị trí của báo cáo so với một số nghiên cứu liên quan.}
\label{tab:related-work}
\small\sloppy
\setlength{\tabcolsep}{3pt}
\begin{tabular}{L{3.0cm}L{4.2cm}L{7.4cm}}
\toprule
Nguồn & Nội dung & Liên hệ với báo cáo \\\midrule
Saxena và cộng sự~\cite{saxena2008damage} & Mô phỏng suy giảm, tạo C-MAPSS &
Cơ sở để hiểu dữ liệu và cách tạo nhãn RUL. \\
Ramasso và Saxena~\cite{ramasso2014review} & Tổng quan các phương pháp trên C-MAPSS &
Lý do báo cáo mô tả rõ subset, cách chia dữ liệu và metric. \\
Lei và cộng sự~\cite{lei2018machinery} & Tổng quan machinery prognostics &
Bối cảnh chung từ thu thập dữ liệu đến dự đoán RUL. \\
Li và cộng sự~\cite{li2018rul} & CNN trên window thời gian &
Tham khảo độ dài window 30 cycles và NASA score. \\
Zheng và cộng sự~\cite{zheng2017lstm} & LSTM cho dự đoán RUL &
Hướng deep learning chưa được thực hiện trong báo cáo. \\
Saxena và cộng sự~\cite{saxena2010metrics} & Metric đánh giá prognostics &
Phân biệt đánh giá tại một điểm cắt với đánh giá trên toàn chuỗi. \\
\bottomrule
\end{tabular}
\end{table}

\subsection{Metric đánh giá}
Với $n$ động cơ, sai số $e_i=\hat{y}_i-y_i$ là RUL dự đoán trừ RUL thật. MAE
bằng $\frac{1}{n}\sum|e_i|$ và RMSE bằng $\sqrt{\frac{1}{n}\sum e_i^2}$; cả hai
có đơn vị cycle, trong đó RMSE nhạy hơn với các sai số lớn. NASA score cộng
điểm phạt của từng dự đoán~\cite{li2018rul}:
\begin{equation}
 S=\sum_{i=1}^{n}s(e_i),\qquad
 s(e)=
 \begin{cases}
  \exp(-e/13)-1, & e<0,\\
  \exp(e/10)-1, & e\geq 0.
 \end{cases}
 \label{eq:nasa-score}
\end{equation}
Mẫu số 10 nhỏ hơn 13, nên dự đoán cao hơn RUL thật (báo hỏng muộn) bị phạt nặng
hơn dự đoán thấp hơn cùng độ lớn, phản ánh việc báo muộn có thể khiến động cơ
hỏng trước khi kịp bảo trì~\cite{saxena2008damage}. NASA score không có đơn vị và tăng theo hàm
mũ, nên chỉ dùng để so sánh các mô hình trên cùng tập đánh giá.

Theo Saxena và cộng sự~\cite{saxena2010metrics}, một dự đoán tại cycle cuối của
mỗi động cơ test chỉ đo chất lượng tại điểm cắt đó, không đo sai số dọc theo
toàn bộ vòng đời.

\subsection{Trailing window và đặc trưng thống kê}
Tại cycle $t$, mẫu đầu vào của động cơ $i$ là ma trận gồm $w$ cycles gần nhất:
\begin{equation}
 W_{i,t}=[\mathbf{x}_{i,t-w+1},\ldots,\mathbf{x}_{i,t}]^{\mathsf T}
 \in\mathbb{R}^{w\times J},\qquad t\geq w,
 \label{eq:window}
\end{equation}
trong đó $\mathbf{x}_{i,t}$ là vector $J=21$ giá trị sensor tại cycle $t$. Báo
cáo dùng $w=30$, một độ dài từng được sử dụng cho FD001~\cite{li2018rul}. Window
chỉ chứa quá khứ và hiện tại, nên không có thông tin từ cycle tương lai; ví dụ,
mẫu tại cycle 170 dùng cycles 141--170.

Mỗi cột sensor trong window được tóm tắt bằng sáu thống kê: giá trị cuối, trung
bình, độ lệch chuẩn mẫu, min, max và độ dốc tuyến tính. Với dãy $z_1,\ldots,z_w$,
độ dốc là hệ số của hồi quy bình phương tối thiểu theo vị trí $k$~\cite{hyndman2021fpp3}:
\begin{equation}
 b=\frac{\sum_{k=1}^{w}(k-\bar{k})(z_k-\bar{z})}
 {\sum_{k=1}^{w}(k-\bar{k})^2},
 \label{eq:slope}
\end{equation}
với $\bar{k}$ và $\bar{z}$ là trung bình của vị trí và của giá trị sensor. Đơn
vị của $b$ là đơn vị sensor trên một cycle. Giá trị cuối và trung bình cho biết
mức tín hiệu, còn độ dốc cho biết tín hiệu đang tăng hay giảm trong 30 cycles
gần nhất.

\subsection{Kiểm định ghép cặp và hiệu chỉnh đa kiểm định}
Khi so sánh một sensor ở đầu và cuối vòng đời của cùng một động cơ, phép ghép
cặp loại bỏ khác biệt nền giữa các động cơ. Kiểm định dấu-hạng Wilcoxon xếp hạng
độ lớn các chênh lệch rồi so sánh tổng hạng của chênh lệch dương và âm; dưới giả
định phân phối chênh lệch đối xứng, nó kiểm định trung vị chênh lệch bằng
0~\cite{wilcoxon1945,scipyWilcoxon}. Bootstrap lấy mẫu lại có hoàn lại các động
cơ để ước lượng khoảng tin cậy cho trung bình chênh lệch mà không cần giả định
phân phối chuẩn~\cite{efron1979bootstrap,scipyBootstrap}.

Khi kiểm định đồng thời 21 sensor, xác suất có ít nhất một kết quả dương tính giả
tăng lên. Thủ tục Benjamini--Hochberg điều chỉnh các p-value để kiểm soát tỷ lệ
phát hiện sai kỳ vọng (FDR) trong nhóm kiểm định~\cite{benjamini1995fdr,scipyFDR}.

\subsection{Health Index và ARIMA}
Health Index là một chỉ số duy nhất tóm tắt trạng thái của động cơ từ nhiều
sensor. Báo cáo tạo Health Index bằng cách chuẩn hóa 21 sensor rồi lấy thành
phần chính thứ nhất (PC1) của PCA~\cite{sklearnPCA}. Dấu của PC1 được chọn sao
cho chỉ số giảm dần theo cycle, tức giá trị thấp ứng với trạng thái xấu hơn.

Chuỗi Health Index $h_t$ của từng động cơ được mô hình hóa bằng
ARIMA$(p,d,q)$~\cite{hyndman2021fpp3}:
\begin{equation}
\phi(B)(1-B)^d h_t = c + \theta(B)\varepsilon_t,
\label{eq:arima}
\end{equation}
trong đó $B$ là toán tử trễ, $\phi(B)$ và $\theta(B)$ là đa thức tự hồi quy bậc
$p$ và trung bình trượt bậc $q$, $d$ là bậc sai phân và $\varepsilon_t$ là nhiễu
trắng. Kiểm định ADF đánh giá tính dừng của chuỗi~\cite{statsmodelsADF}, còn
kiểm định Ljung--Box kiểm tra phần dư sau khi fit còn tự tương quan hay không.

\subsection{Xử lý dữ liệu với Spark và lưu trữ phân lớp}
Apache Spark xử lý dữ liệu phân tán dựa trên Resilient Distributed Dataset,
cho phép tính toán song song và phục hồi khi lỗi~\cite{zaharia2012rdd}; API
DataFrame của Spark SQL hỗ trợ các phép biến đổi trên dữ liệu có
schema~\cite{sparkSQLGuide}. Tổ chức dữ liệu theo ba lớp Bronze--Silver--Gold
tách dữ liệu gốc, dữ liệu đã làm sạch và bảng tổng hợp, giúp mỗi bước biến đổi
có đầu vào và đầu ra rõ ràng. Parquet lưu dữ liệu theo cột, phù hợp cho truy vấn
phân tích chỉ đọc một số cột.
